\documentclass[tikz,border=0pt]{standalone}
% Shared publication style. Compile figure sources from this directory.
\usepackage[T1]{fontenc}
\usepackage{lmodern}
\usepackage{amsmath,amssymb}
\usepackage{tikz}
\usepackage{pgfplots}
\usetikzlibrary{arrows.meta,calc,positioning,backgrounds,patterns,shapes.geometric}
\usepgfplotslibrary{groupplots,fillbetween}
\pgfplotsset{compat=1.18}
\definecolor{fgblue}{HTML}{0072B2}
\definecolor{fgorange}{HTML}{D55E00}
\definecolor{fgslate}{HTML}{8B949E}
\definecolor{fgdark}{HTML}{25313C}
\definecolor{fgpale}{HTML}{F2F6F8}
\definecolor{fggrid}{HTML}{E3E8EC}
\definecolor{fgwater}{HTML}{D6EAF4}
\definecolor{fgland}{HTML}{E8E3D3}
\renewcommand{\familydefault}{\sfdefault}
\tikzset{every picture/.style={font=\fontsize{9}{11}\selectfont,text=fgdark},
  panel label/.style={font=\bfseries\fontsize{10}{12}\selectfont,anchor=north west},
  panel title/.style={font=\bfseries\fontsize{9}{11}\selectfont,anchor=north west},
  light arrow/.style={-{Stealth[length=1.6mm]},draw=fgslate,line width=.65pt}}
\pgfplotsset{fg axis/.style={
  axis line style={fgslate,line width=.4pt},axis x line*=bottom,axis y line*=left,
  tick align=outside,tick style={fgslate,line width=.4pt},
  tick label style={font=\fontsize{8}{10}\selectfont,text=fgdark},
  label style={font=\fontsize{9}{11}\selectfont,text=fgdark},
  title style={font=\bfseries\fontsize{9}{11}\selectfont,text=fgdark},
  grid style={fggrid,line width=.35pt},minor tick num=0,
  legend style={draw=none,fill=none,font=\fontsize{8}{10}\selectfont},
  scaled ticks=false,clip=false,
  every axis plot/.append style={line width=.85pt,mark size=2.2pt}
}}

% B--C use the SAME regenerated initial seed-11 environment draw.
% See make_tail_data.py and fig02_data_provenance.json.
\begin{document}
\begin{tikzpicture}[x=1mm,y=1mm]
\path[use as bounding box] (0,0) rectangle (140,61);
\node[panel label] at (0,61) {A};
\node[panel title] at (5,61) {Score geometry};
\node[panel label] at (44,61) {B};
\node[panel title] at (49,61) {Information collapse};
\node[panel label] at (91,61) {C};
\node[panel title] at (96,61) {Inverse-information tail};

% The drawing is a normalized L2 score-space illustration, not data.
\fill[fgslate!10] (1,34.8) rectangle (35,36.2);
\draw[fgslate,line width=.55pt] (1,35.5)--(35,35.5);
\draw[fgslate,line width=1.4pt,-{Stealth[length=1.8mm]}] (3,35.5)--(29,35.5);
\draw[fgdark,line width=1pt,-{Stealth[length=2mm]}] (3,35.5)--(29,42);
\draw[fgblue,line width=1.5pt,-{Stealth[length=1.6mm]}] (29,35.5)--(29,42);
\draw[fgslate,line width=.4pt] (27.5,35.5)--(27.5,37)--(29,37);
\fill[fgdark] (3,35.5) circle (.65);
\node[anchor=south,font=\fontsize{8}{10}\selectfont] at (15,41.8) {Task score};
\node[anchor=north,font=\fontsize{8}{10}\selectfont,text=fgslate!70!black] at (17.5,33.4) {Nuisance span};
\draw[fgblue!65,line width=.45pt] (30.7,39)--(34,39)--(34,23.5)--(29,23.5);
\node[anchor=east,align=right,text=fgblue,font=\fontsize{8}{10}\selectfont] at (28.7,23.5) {Remaining\\task score};
\node[font=\fontsize{12}{14}\selectfont,text=fgblue] at (18,10) {$\lambda=\delta^2$};

\begin{axis}[fg axis,at={(51mm,15mm)},anchor=south west,width=33mm,height=35mm,scale only axis,
 xmin=0,xmax=1,ymode=log,ymin=1e-9,ymax=3,ytick={1e-8,1e-4,1},clip=true,
 xtick={0,.5,1},xticklabels={0,.5,1},grid=major,
 xlabel={Environment quantile},xlabel style={font=\fontsize{8}{10}\selectfont},ylabel={Information},
 ylabel style={font=\fontsize{8}{10}\selectfont},log basis y=10]
\addplot[fgslate,line width=1pt] table[x=quantile,y=raw] {fig02_quantiles.dat};
\addplot[fgblue,line width=1.2pt] table[x=quantile,y=projected] {fig02_quantiles.dat};
\node[anchor=north,text=fgslate!70!black,font=\fontsize{8}{10}\selectfont] at (axis cs:.71,.9) {Raw};
\node[anchor=north,text=fgblue,font=\fontsize{8}{10}\selectfont] at (axis cs:.7,.002) {Remaining};
\end{axis}
\begin{axis}[fg axis,at={(103mm,15mm)},anchor=south west,width=35mm,height=35mm,scale only axis,
 xmode=log,ymode=log,xmin=30,xmax=5e7,ymin=3e-4,ymax=1,clip=true,
 xtick={1e2,1e4,1e6},ytick={1,.01,.001},grid=major,
 xlabel={Inverse information $x$},xlabel style={font=\fontsize{8}{10}\selectfont},
 ylabel={$\Pr(L>x)$},ylabel style={font=\fontsize{8}{10}\selectfont}]
\addplot[fgslate,dashed,line width=.8pt] table[x=threshold,y=theory] {fig02_survival.dat};
\addplot[only marks,mark=*,mark size=1pt,fgblue,mark repeat=10] table[x=threshold,y=survival] {fig02_survival.dat};
\node[anchor=south west,font=\fontsize{8}{10}\selectfont] at (axis cs:1500,.12) {Slope $-1/2$};
\end{axis}
\end{tikzpicture}
\end{document}
